% !TeX root = ../../../main.tex
\section{Torsors and Rank-One Local Systems}
% $X$ with etale site, but say things here work for toposes in general and flat sites in general.

In what follows, we largely adhere to the treatment of torsors and group objects found in the expository notes of Alex Youcis \cite{YoucisTorsors}.
Let $\sC=(\cC, J)$ be a site and let $\cE=Sh(\sC)$ be the associated topos. 
Let $\mathcal{G}$ be a group object in $\mathcal{E}$. 
We denote by $\mathcal{G}\mathcal{E}$ the category of objects in $\mathcal{E}$ endowed with a left $\mathcal{G}$-action. 

\begin{definition}
    A \textbf{$\mathcal{G}$-torsor} in $\mathcal{E}$ is an object $\mathcal{P}$ of $\mathcal{G}\mathcal{E}$ satisfying the following conditions:
    \begin{enumerate}
        \item The structural morphism $\mathcal{P} \to 1$ is an epimorphism in $\mathcal{E}$ (i.e., $\mathcal{P}$ is locally non-empty).
        \item The map $\mathcal{G} \times \mathcal{P} \to \mathcal{P} \times \mathcal{P}$ defined by $(g, p) \mapsto (g \cdot p, p)$ is an isomorphism in $\mathcal{E}$ (i.e., $\mathcal{G}$ acts simply transitively on $\mathcal{P}$).
    \end{enumerate}
\end{definition}

Since $\mathcal{E}$ is the topos of sheaves on a site $\sC$, the definition can be reformulated in terms of covers. A $\mathcal{G}$-sheaf $\mathcal{P}$ is a $\mathcal{G}$-torsor if:
\begin{enumerate}
    \item For every object $X \in \mathcal{C}$, there exists a covering $\{ U_i \to X \} \in J$ such that $\mathcal{P}(U_i) \neq \emptyset$ for all $i$.
    \item For any $X \in \mathcal{C}$ where $\mathcal{P}(X)$ is non-empty, the action of $\mathcal{G}(X)$ on $\mathcal{P}(X)$ is simply transitively.
\end{enumerate}

A fundamental property of torsors is their local triviality: 
a $\mathcal{G}$-sheaf $\mathcal{P}$ is a $\mathcal{G}$-torsor if and only if it is locally isomorphic to the trivial torsor. 
Specifically, for every $X \in \mathcal{C}$, there must exist a cover $\{ U_i \to X \}$ such that the restriction $\mathcal{P}|_{U_i}$ is 
isomorphic, as a $\mathcal{G}|_{U_i}$-sheaf, to $\mathcal{G}|_{U_i}$ acting on itself by left multiplication.

A \textbf{morphism of $\mathcal{G}$-torsors} $f: \mathcal{P}_1 \to \mathcal{P}_2$ is a morphism of sheaves that is equivariant with respect to the $\mathcal{G}$-action.

\begin{prop}\label{prop:torsor_morphism_iso}
    Every morphism of $\mathcal{G}$-torsors is an isomorphism.
    Consequently, the category of $\mathcal{G}$-torsors in $\mathcal{E}$ forms a groupoid.
\end{prop}
\begin{proof}
    Locally on a cover trivializing both torsors, such a morphism is a $\mathcal{G}$-equivariant endomorphism
    of the trivial torsor $\mathcal{G}$, i.e.\ right translation by a section of $\mathcal{G}$, hence an
    isomorphism; and being an isomorphism may be checked on a cover.
\end{proof}

\begin{definition}
    We denote the groupoid of $\mathcal{G}$-torsors in $\mathcal{E}$ by $\mathbf{Tors}(\mathcal{E}, \mathcal{G})$ .  The set of isomorphism classes of $\mathcal{G}$-torsors is denoted by $\mathrm{Tors}(\mathcal{E}, \mathcal{G})$.
\end{definition}

Torsors exhibit functoriality with respect to the group object:

\begin{definition}
Let $\varphi: \mathcal{G}_1 \to \mathcal{G}_2$ be a morphism of group 
sheaves on $\sC$, and let $\mathcal{P}$ be a $\mathcal{G}_1$-torsor. We define the \textbf{contracted product} $\mathcal{G}_2 \times^{\mathcal{G}_1} \mathcal{P}$ as the quotient sheaf $\mathcal{G}_1 \backslash (\mathcal{G}_2 \times \mathcal{P})$, where $\mathcal{G}_1$ acts on the product by:
\[ g_1 \cdot (g_2, p) = (g_2 \varphi(g_1)^{-1}, g_1 \cdot p) \]

The contracted product inherits a natural left $\mathcal{G}_2$-action given on local sections by $h \cdot [g_2, p] = [h g_2, p]$, which endows it with the structure of a $\mathcal{G}_2$-torsor. 
\end{definition}

This construction yields a functor:
\[ \varphi_*: \mathbf{Tors}(\cE, \mathcal{G}_1) \to \mathbf{Tors}(\cE, \mathcal{G}_2), \quad \mathcal{P} \mapsto \mathcal{G}_2 \times^{\mathcal{G}_1} \mathcal{P} \]
There is a canonical $\varphi$-equivariant morphism
\[
j_\varphi:\mathcal P\longrightarrow\varphi_*\mathcal P,
\qquad p\longmapsto[1,p],
\]
meaning that
$j_\varphi(g\cdot p)=\varphi(g)\cdot j_\varphi(p)$ for every local section
$g$ of $\mathcal G_1$.
On the level of isomorphism classes, $\varphi_*$ induces a map of pointed sets $\mathrm{Tors}(\cE, \mathcal{G}_1) \to \mathrm{Tors}(\cE, \mathcal{G}_2)$, sending the class of the trivial $\mathcal{G}_1$-torsor to the class of the trivial $\mathcal{G}_2$-torsor.

When $\mathcal{G}$ is a \textbf{sheaf of abelian groups} (an \textit{abelian sheaf}), the pointed set $\mathrm{Tors}(\cE, \mathcal{G})$ 
inherits the structure of an abelian group:
\begin{definition}
    Assume $\mathcal{G}$ is abelian sheaf. Let $\mathcal{P}_1$ and $\mathcal{P}_2$ be objects of $\mathbf{Tors}(\cE, \mathcal{G})$. 
    We define the \textbf{product} $\mathcal{P}_1 \otimes^{\cG} \mathcal{P}_2$ to be the quotient sheaf 
    $(\mathcal{P}_1 \times \mathcal{P}_2)/\mathcal{G}$, where $\mathcal{G}$ acts on $T$-points by having:
    \[ g \cdot (f_1, f_2) := (g f_1, g^{-1} f_2) \]
    The group object $\mathcal{G}$ then acts on the resulting quotient via its action on the presheaf quotient, which is given on classes by:
    \[ g \cdot [(f_1, f_2)] = [(g f_1, f_2)] = [(f_1, g f_2)] \]
    where the square brackets denote the class in the quotient set.\footnote{
    Equivantly, the sum is obtained as the contracted product of the $\cG \times \cG$-torsor $\cP_1 \times \cP_2$ along the multiplication map
$m: \cG \times \cG \to \cG$. }
\end{definition}

Let $[\mathcal{P}_1]$ and $[\mathcal{P}_2]$ be the isomorphism classes of $\mathcal{P}_1$ and $\mathcal{P}_2$ in $\mathrm{Tors}(\cE, \mathcal{G})$.
We define the sum $[\mathcal{P}_1] + [\mathcal{P}_2]$ to be the class $[\cP_1 \otimes^{\cG} \cP_2]$.
This structure turns $\mathrm{Tors}(\cE, \mathcal{G})$ into an abelian group, 
where the identity is the class of the trivial torsor and the inverse is obtained by the opposite action.

We now record the equivalence used in the thesis.
Let $\Lambda$ be a commutative ring object of $\mathcal{E}$ and let
$G=\Lambda^\times$ be its group of units. We write
$\mathbf{Pic}(\mathcal{E},\Lambda)$ for the groupoid of locally free
$\Lambda$-modules of rank one. For $\mathcal{L}$ in this groupoid, its
\textbf{frame torsor} is
\[
\operatorname{Fr}(\mathcal{L})
  :=\underline{\operatorname{Isom}}_\Lambda(\Lambda,\mathcal{L}).
\]
The group $G$ acts on frames by precomposition with multiplication on $\Lambda$.
Conversely, if $\mathcal{P}$ is a $G$-torsor, its \textbf{associated
$\Lambda$-module} is the contracted product
\[
\mathcal{P}\times^G\Lambda
  :=(\mathcal{P}\times\Lambda)/G,
\qquad
g\cdot(p,\lambda)=(gp,g^{-1}\lambda),
\]
where $G$ acts on $\Lambda$ by multiplication.

\begin{prop}
\label{prop:torsor_module_equivalence}
The functors
\[
\mathcal{P}\longmapsto\mathcal{P}\times^G\Lambda,
\qquad
\mathcal{L}\longmapsto\operatorname{Fr}(\mathcal{L})
\]
are mutually quasi-inverse equivalences
\[
\mathbf{Tors}(\mathcal{E},G)
\xrightarrow{\ \sim\ }
\mathbf{Pic}(\mathcal{E},\Lambda).
\]
\end{prop}

\begin{proof}
There are natural morphisms
\[
\mathcal{P}\longrightarrow
\operatorname{Fr}(\mathcal{P}\times^G\Lambda),
\qquad
p\longmapsto\bigl(\lambda\longmapsto[p,\lambda]\bigr),
\]
and
\[
\operatorname{Fr}(\mathcal{L})\times^G\Lambda
\longrightarrow\mathcal{L},
\qquad
[\varphi,\lambda]\longmapsto\varphi(\lambda).
\]
After passing to a cover on which $\mathcal{P}$ is the trivial torsor and
$\mathcal{L}$ is free of rank one, both morphisms become the identity.
Hence they are isomorphisms.
\end{proof}

When $\mathcal{E}=\operatorname{Sh}(X_{\text{\'et}})$ and $\Lambda$
is a finite constant ring, the objects of
$\mathbf{Pic}(\mathcal{E},\Lambda)$ are precisely the rank-one
$\Lambda$-local systems on $X$. Henceforth, we identify such a local system
with its $G=\Lambda^\times$-torsor of frames and pass between the two
descriptions without further comment.\footnote{The equivalence is compatible
with pullback and with tensor products, hence with the symmetric-power
constructions used below. These compatibilities may be checked after a local
trivialization.}

Lastly, for any object $\cX \in \mathcal{E}$, there is a canonical identification between the slice category
$(\mathcal{G}\mathcal{E})/\cX$ and the category of group objects $\mathcal{G}(\mathcal{E}/\cX)$, where $\cX$ is viewed as having the 
trivial $\mathcal{G}$-action.

We denote by $\mathbf{Tors}(\cX, \cG)$ the category of $G$-torsors over $\cX$ in $\cG \cE_{/\cX}$.

%%% %%% %%% %%% %%% %%% %%% %%%
%       Torsors as Spaces
%%% %%% %%% %%% %%% %%% %%% %%%
